\documentclass[10pt,a4paper]{amsart}
\usepackage{amsmath,amssymb,amsthm,graphicx}
% Retained minimal preamble for the Stacks Project chapter more-algebra.tex (More on Algebra).
% The chapter file pinned-original-source.tex carries only \input{preamble} on its first line, so no
% preamble text lives inside the chapter itself.  The Stacks Project uses one single shared preamble for
% every chapter; the snapshot stages it identically as sources/stacks-chow/preamble.tex and
% sources/stacks-etale-cohomology/preamble.tex.  The block below is copied mechanically, line for line,
% from sources/stacks-chow/preamble.tex lines 14-18 (the Xy-pic setup) and 154-217 (the native theorem
% declarations together with the mathematical macros).  Nothing was added, removed or redefined; the
% wrapper class and the amsmath/amssymb/amsthm/graphicx base packages are supplied by the assembler.
% Chapter: more-algebra.tex -- More on Algebra.
% For commutative diagrams we use Xy-pic
\usepackage[all]{xy}

% We use 2cell for 2-commutative diagrams.
\xyoption{2cell}

% Theorem environments.
%
\theoremstyle{plain}
\newtheorem{theorem}[subsection]{Theorem}
\newtheorem{proposition}[subsection]{Proposition}
\newtheorem{lemma}[subsection]{Lemma}

\theoremstyle{definition}
\newtheorem{definition}[subsection]{Definition}
\newtheorem{example}[subsection]{Example}
\newtheorem{exercise}[subsection]{Exercise}
\newtheorem{situation}[subsection]{Situation}

\theoremstyle{remark}
\newtheorem{remark}[subsection]{Remark}
\newtheorem{remarks}[subsection]{Remarks}

\numberwithin{equation}{subsection}

% Macros
%
\def\lim{\mathop{\mathrm{lim}}\nolimits}
\def\colim{\mathop{\mathrm{colim}}\nolimits}
\def\Spec{\mathop{\mathrm{Spec}}}
\def\Hom{\mathop{\mathrm{Hom}}\nolimits}
\def\Ext{\mathop{\mathrm{Ext}}\nolimits}
\def\SheafHom{\mathop{\mathcal{H}\!\mathit{om}}\nolimits}
\def\SheafExt{\mathop{\mathcal{E}\!\mathit{xt}}\nolimits}
\def\Sch{\mathit{Sch}}
\def\Mor{\mathop{\mathrm{Mor}}\nolimits}
\def\Ob{\mathop{\mathrm{Ob}}\nolimits}
\def\Sh{\mathop{\mathit{Sh}}\nolimits}
\def\NL{\mathop{N\!L}\nolimits}
\def\CH{\mathop{\mathrm{CH}}\nolimits}
\def\proetale{{pro\text{-}\acute{e}tale}}
\def\etale{{\acute{e}tale}}
\def\QCoh{\mathit{QCoh}}
\def\Ker{\mathop{\mathrm{Ker}}}
\def\Im{\mathop{\mathrm{Im}}}
\def\Coker{\mathop{\mathrm{Coker}}}
\def\Coim{\mathop{\mathrm{Coim}}}

% Boxtimes
%
\DeclareMathSymbol{\boxtimes}{\mathbin}{AMSa}{"02}

%
% Macros for moduli stacks/spaces
%
\def\QCohstack{\mathcal{QC}\!\mathit{oh}}
\def\Cohstack{\mathcal{C}\!\mathit{oh}}
\def\Spacesstack{\mathcal{S}\!\mathit{paces}}
\def\Quotfunctor{\mathrm{Quot}}
\def\Hilbfunctor{\mathrm{Hilb}}
\def\Curvesstack{\mathcal{C}\!\mathit{urves}}
\def\Polarizedstack{\mathcal{P}\!\mathit{olarized}}
\def\Complexesstack{\mathcal{C}\!\mathit{omplexes}}
% \Pic is the operator that assigns to X its picard group, usage \Pic(X)
% \Picardstack_{X/B} denotes the Picard stack of X over B
% \Picardfunctor_{X/B} denotes the Picard functor of X over B
\def\Pic{\mathop{\mathrm{Pic}}\nolimits}
\def\Picardstack{\mathcal{P}\!\mathit{ic}}
\def\Picardfunctor{\mathrm{Pic}}
\def\Deformationcategory{\mathcal{D}\!\mathit{ef}}

% --- appended 2026-10-12 (Chief): the project own definitions of the reference/slogan/history
% display environments, copied mechanically from the shared Stacks preamble lines 8-12
% (\usepackage{verbatim} plus the three \newenvironment lines, which the chapter preamble keeps
% verbatim from the project preamble). Nothing added, removed or redefined: these environments
% carry no mathematical content and the project itself typesets them as comments. Without these
% three lines a counted statement whose own text uses them fails with Environment slogan undefined.
\usepackage{verbatim}
\newenvironment{reference}{\comment}{\endcomment}
\newenvironment{slogan}{\comment}{\endcomment}
\newenvironment{history}{\comment}{\endcomment}
\usepackage{cleveref}
\crefname{originallemma}{original lemma}{original lemmas}
\Crefname{originallemma}{Original lemma}{Original lemmas}
\makeatletter
\newlabel{more-algebra-lemma-formally-smooth}{{07EC}{source}{}{sourcepointer-1}{}}
\newlabel{more-algebra-lemma-formally-smooth@cref}{{[originallemma][0][]07EC}{[1][0][]source}{}{}{}}
\makeatother
\crefname{originallemma}{original lemma}{original lemmas}
\Crefname{originallemma}{Original lemma}{Original lemmas}
\makeatletter
\newlabel{algebra-definition-formally-smooth}{{00TI}{source}{}{sourcepointer-2}{}}
\newlabel{algebra-definition-formally-smooth@cref}{{[originallemma][0][]00TI}{[1][0][]source}{}{}{}}
\makeatother
\crefname{originallemma}{original lemma}{original lemmas}
\Crefname{originallemma}{Original lemma}{Original lemmas}
\makeatletter
\newlabel{definition-formally-smooth}{{07EB}{source}{}{sourcepointer-3}{}}
\newlabel{definition-formally-smooth@cref}{{[originallemma][0][]07EB}{[1][0][]source}{}{}{}}
\makeatother
\crefname{originallemma}{original lemma}{original lemmas}
\Crefname{originallemma}{Original lemma}{Original lemmas}
\makeatletter
\newlabel{lemma-continuous}{{07E9}{source}{}{sourcepointer-4}{}}
\newlabel{lemma-continuous@cref}{{[originallemma][0][]07E9}{[1][0][]source}{}{}{}}
\makeatother
\title{Formal smoothness for adic topologies follows from formal smoothness, and the discrete variant}
\author{The Stacks Project Authors}
\date{stacks-project:more-algebra.tex:local-snapshot-20261012; GNU FDL 1.2 or later}
\begin{document}
\maketitle
\setcounter{subsection}{1}
\noindent\textit{Editorial note.} All delivered mathematical text is the author's own, copied from the pinned Stacks Project chapter file more-algebra.tex; no statement or proof step was rewritten, reordered or removed.\par\medskip
\noindent\textit{Editorial note.} The chapter file pinned-original-source.tex carries only \textbackslash{}input\{preamble\} on its first line, so no preamble text lives inside the chapter. The retained file sources/stacks-more-algebra/retained-preamble.tex is a mechanical, line-for-line copy of the project's single shared preamble as staged in the snapshot (Xy-pic setup plus theorem declarations and mathematical macros). Nothing was added, removed or redefined, and the excerpt serves only as preamble, never as a span coordinate.\par\medskip
\noindent\textit{Editorial note.} Both counted spans start at the \textbackslash{}begin line of their environment, so the statement and the author's proof are delivered whole; surrounding prose of the section is deliberately outside the counted spans and is carried as standard prerequisites instead.\par\medskip
\noindent\textit{Editorial note.} The citation fields (source\_id, source\_version, authors, source\_url, license.name, source\_title) are identical to those of every other card this lane has produced from this chapter, so the batch selector treats them as one source.\par\medskip
\noindent\textit{Allowed standard prerequisites (editorial).}\begin{itemize}
\item Formal smoothness in the algebraic sense and its defining lifting property. Algebra, Definition \textbackslash{}ref\{algebra-definition-formally-smooth\} and Definition \textbackslash{}ref\{definition-formally-smooth\}, cited at proof lines 9186-9188 and 9194. Supply the lifting property that is checked against the topological continuity requirements.
\item A pre-adic topology on a ring is the adic topology of some ideal. used at proof lines 9188-9190. Lets the author name the ideal defining the topology on the target.
\item Continuity of a ring map for adic topologies is equivalent to a power of the source ideal mapping into the target ideal. Lemma \textbackslash{}ref\{lemma-continuous\}, cited at proof lines 9191 and 9199-9200 and 9206-9208. Translates the continuity of the structure maps and of the constructed lifts into statements about powers of ideals.
\end{itemize}\medskip
\ifdefined\hypertarget\hypertarget{sourcepointer-1}{}\fi
\noindent\textit{Source reference.} Original lemma 07EC: focused-seeds/AGP/sources/stacks-snapshot-20261011/tags.txt, source lines 9442--9442. The lemma's own canonical tag, quoted so that readers can resolve the label inside the delivered text; it points at the counted statement itself and carries no proof content.\par\medskip
\ifdefined\hypertarget\hypertarget{sourcepointer-2}{}\fi
\noindent\textit{Source reference.} Original lemma 00TI: focused-seeds/AGP/sources/stacks-snapshot-20261011/tags.txt, source lines 1076--1076. Navigation pointer to a statement the counted proof refers to but does not contain; nothing load-bearing is delegated to it.\par\medskip
\ifdefined\hypertarget\hypertarget{sourcepointer-3}{}\fi
\noindent\textit{Source reference.} Original lemma 07EB: focused-seeds/AGP/sources/stacks-snapshot-20261011/tags.txt, source lines 9441--9441. Navigation pointer to a statement the counted proof refers to but does not contain; nothing load-bearing is delegated to it.\par\medskip
\ifdefined\hypertarget\hypertarget{sourcepointer-4}{}\fi
\noindent\textit{Source reference.} Original lemma 07E9: focused-seeds/AGP/sources/stacks-snapshot-20261011/tags.txt, source lines 9439--9439. Navigation pointer to a statement the counted proof refers to but does not contain; nothing load-bearing is delegated to it.\par\medskip
\begin{lemma}
\label{lemma-formally-smooth}
Let $\varphi : R \to S$ be a ring map.
\begin{enumerate}
\item If $R \to S$ is formally smooth in
the sense of Algebra, Definition \ref{algebra-definition-formally-smooth},
then $R \to S$ is formally smooth for any linear topology on $R$ and
any pre-adic topology on $S$ such that $R \to S$ is continuous.
\item Let $\mathfrak n \subset S$ and $\mathfrak m \subset R$
ideals such that $\varphi$ is continuous for the $\mathfrak m$-adic
topology on $R$ and the $\mathfrak n$-adic topology
on $S$. Then the following are equivalent
\begin{enumerate}
\item $\varphi$ is formally smooth for the $\mathfrak m$-adic topology on
$R$ and the $\mathfrak n$-adic topology on $S$, and
\item $\varphi$ is formally smooth for the discrete topology
on $R$ and the $\mathfrak n$-adic topology on $S$.
\end{enumerate}
\end{enumerate}
\end{lemma}

\begin{proof}
Assume $R \to S$ is formally smooth in
the sense of Algebra, Definition \ref{algebra-definition-formally-smooth}.
If $S$ has a pre-adic topology, then
there exists an ideal $\mathfrak n \subset S$ such that $S$ has the
$\mathfrak n$-adic topology. Suppose given a solid commutative diagram as in
Definition \ref{definition-formally-smooth}.
Continuity of $S \to A/J$ means that $\mathfrak n^k$ maps to zero
in $A/J$ for some $k \geq 1$, see Lemma \ref{lemma-continuous}.
We obtain a ring map $\psi : S \to A$ from the assumed formal smoothness of
$S$ over $R$. Then $\psi(\mathfrak n^k) \subset J$ hence
$\psi(\mathfrak n^{2k}) = 0$ as $J^2 = 0$. Hence $\psi$ is continuous by
Lemma \ref{lemma-continuous}. This proves (1).

\medskip\noindent
The proof of (2)(b) $\Rightarrow$ (2)(a) is the same as the proof of (1).
Assume (2)(a). Suppose given a solid commutative diagram as in
Definition \ref{definition-formally-smooth} where we use the discrete
topology on $R$. Since $\varphi$ is continuous we see that
$\varphi(\mathfrak m^n) \subset \mathfrak n$ for some $n \geq 1$.
As $S \to A/J$ is continuous we see that $\mathfrak n^k$ maps to
zero in $A/J$ for some $k \geq 1$. Hence $\mathfrak m^{nk}$ maps
into $J$ under the map $R \to A$. Thus $\mathfrak m^{2nk}$ maps to zero
in $A$ and we see that $R \to A$ is continuous in the $\mathfrak m$-adic
topology. Thus (2)(a) gives a dotted arrow as desired.
\end{proof}

\end{document}
